% =============================================================
%  05_strutture_dati.tex
%  Capitolo 5 --- Strutture dati
% =============================================================

\chapter{Strutture dati}

\begin{intuizione}
Una lista è flessibile ma lenta per l'accesso per chiave.
Una hash table è veloce ma non ordinata.
\fn{defstruct} dà nomi ai campi.
Scegliere la struttura giusta non è un dettaglio: determina
la complessità dell'algoritmo.
\end{intuizione}

\vspace{4pt}

% ================================================================
\section{Association list (alist)}
% ================================================================

\begin{concetto}{Alist --- la mappa più semplice}
\keyline{}{Una alist è una lista di coppie \fn{(chiave . valore)}.}
\keyline{\fn{assoc}}{trova la prima coppia con una data chiave. $O(n)$.}
\keyline{\fn{acons}}{aggiunge una coppia in testa. $O(1)$.}

\boxsep
\detail{La alist è la struttura associativa di base di Lisp.
Funziona bene per dizionari piccoli (< 20--30 coppie) o quando
l'ordine di inserimento conta.
Per dizionari grandi, preferire le hash table.}
\end{concetto}

\begin{esempio}[Alist: costruzione e accesso]
\begin{lstlisting}
;; Costruzione
(defparameter *persona*
  '((nome . "Mario") (eta . 23) (citta . "Roma")))

;; assoc: cerca per chiave
(assoc 'eta *persona*)           ; => (ETA . 23)
(cdr (assoc 'eta *persona*))     ; => 23   (solo il valore)
(assoc 'telefono *persona*)      ; => NIL  (non trovato)

;; acons: aggiunge in testa (non modifica l'originale)
(acons 'email "mario@example.com" *persona*)
; => ((EMAIL . "mario@example.com") (NOME . "Mario") ...)

;; Aggiornare un valore: aggiungi in testa, la vecchia coppia
;; rimane ma non viene trovata (shadow)
(defparameter *p2*
  (acons 'eta 24 *persona*))
(cdr (assoc 'eta *p2*))   ; => 24  (trovata la nuova prima)
\end{lstlisting}
\end{esempio}

\begin{confronto}
\textbf{Alist vs hash table vs \fn{defstruct}.}
\begin{itemize}
  \item \textbf{Alist}: chiavi arbitrarie, funzionale (no mutazione),
        ricerca $O(n)$. Usa per configurazioni, piccoli dizionari,
        risultati di \fn{raggruppa-per}.
  \item \textbf{Hash table}: ricerca $O(1)$, mutabile, chiavi
        confrontate con \fn{equal}/\fn{eql}. Usa per dizionari grandi
        o frequentemente consultati.
  \item \textbf{\fn{defstruct}}: campi nominati, accesso $O(1)$,
        type-checking automatico. Usa quando la struttura è fissa
        e i campi sono noti a priori.
\end{itemize}
\end{confronto}

\begin{workbook}
\textbf{Esercizi 93--94, 100}
(Lookup/set su alist, Frequenza con alist, Conversione alist$\leftrightarrow$hash)\\[4pt]
L'esercizio~93 è il fondamentale: \fn{assoc} e la costruzione
funzionale con \fn{acons}.
L'esercizio~94 (frequenza elementi) mostra come una alist si presta
bene per accumulare conteggi su insiemi piccoli.
\end{workbook}

% ================================================================
\section{Hash table}
% ================================================================

\begin{concetto}{Hash table --- dizionario a tempo costante}
\keyline{\fn{make-hash-table}}{crea una hash table vuota.}
\keyline{\fn{gethash}}{legge un valore per chiave. $O(1)$.}
\keyline{\fn{setf}\ \fn{gethash}}{scrive o sovrascrive un valore. $O(1)$.}
\keyline{\fn{remhash}}{rimuove una coppia. $O(1)$.}
\keyline{\fn{maphash}}{itera su tutte le coppie.}
\end{concetto}

\begin{esempio}[Hash table: operazioni fondamentali]
\begin{lstlisting}
;; Creare
(defparameter *ht* (make-hash-table))

;; Scrivere
(setf (gethash 'nome *ht*) "Mario")
(setf (gethash 'eta  *ht*) 23)

;; Leggere
(gethash 'nome *ht*)    ; => "Mario"  T
(gethash 'tel  *ht*)    ; => NIL      NIL  (due valori: valore + trovato?)

;; Il secondo valore di gethash indica se la chiave esiste
(multiple-value-bind (val trovato)
    (gethash 'eta *ht*)
  (if trovato
      (format t "eta' = ~a~%" val)
      (print "non trovato")))
; => eta' = 23

;; Valore di default esplicito
(gethash 'tel *ht* "sconosciuto")   ; => "sconosciuto"

;; Rimuovere
(remhash 'eta *ht*)
(gethash 'eta *ht*)    ; => NIL  NIL

;; Quante coppie?
(hash-table-count *ht*)   ; => 1
\end{lstlisting}
\end{esempio}

\argomento{Iterare: \fn{maphash}}

\begin{concetto}{\fn{maphash} --- iterare su tutte le coppie}
\begin{sintassi}
(maphash \#'(lambda (chiave valore) ...) ht)
\end{sintassi}
\detail{Non ha un valore di ritorno definito. Usalo per side-effect
(stampare, accumulare in una struttura esterna).}
\end{concetto}

\begin{esempio}[\fn{maphash} e conversione in alist]
\begin{lstlisting}
;; Stampa tutte le coppie
(maphash #'(lambda (k v)
             (format t "~a => ~a~%" k v))
         *ht*)

;; Converti hash table in alist
(defun hash->alist (ht)
  (let ((result '()))
    (maphash #'(lambda (k v)
                 (setf result (acons k v result)))
             ht)
    result))

;; Converti alist in hash table
(defun alist->hash (alist)
  (let ((ht (make-hash-table)))
    (dolist (pair alist ht)
      (setf (gethash (car pair) ht) (cdr pair)))))
\end{lstlisting}
\end{esempio}

\argomento{Test di uguaglianza e stringhe}

\begin{attenzione}
\textbf{Il test di default è \fn{eql}: non funziona con le stringhe.}
\fn{eql} confronta per identità (stesso oggetto in memoria).
Due stringhe con lo stesso contenuto sono due oggetti distinti.

\begin{lstlisting}
;; SBAGLIATO per stringhe: usa :test #'equal
(defparameter *ht-sbagliata* (make-hash-table))
(setf (gethash "foo" *ht-sbagliata*) 42)
(gethash "foo" *ht-sbagliata*)   ; => NIL !! (eql fallisce)

;; CORRETTO: specifica il test di uguaglianza
(defparameter *ht-corretta*
  (make-hash-table :test #'equal))
(setf (gethash "foo" *ht-corretta*) 42)
(gethash "foo" *ht-corretta*)    ; => 42  T
\end{lstlisting}

Regola: chiavi stringa o lista $\to$ \fn{:test \#'equal}.
Chiavi simbolo o numero $\to$ default \fn{eql} va bene.
\end{attenzione}

\begin{workbook}
\textbf{Esercizi 95--100, 116, 120}
(Operazioni base, frequenza alist vs hash, conversioni,
rubrica, conta parole)\\[4pt]
Priorità: \textbf{95, 96, 100}.
L'esercizio~96 è fondamentale: ripete la stessa cosa dell'es.~94
con hash table invece di alist --- il confronto diretto insegna
quando usare quale.\\[4pt]
L'esercizio~120 (conta parole) è il caso reale più comune di hash table:
contatore di frequenze su input testuale di dimensione arbitraria.
\end{workbook}

% ================================================================
\section{Array multidimensionali}
% ================================================================

\begin{concetto}{Array 2D --- matrici}
\keyline{}{Passa una lista di dimensioni a \fn{make-array}.}
\keyline{\fn{aref}}{accede con indici multipli.}
\end{concetto}

\begin{esempio}[Matrice 2D]
\begin{lstlisting}
;; Matrice 3x3 inizializzata a 0
(defparameter *m*
  (make-array '(3 3) :initial-element 0))

;; Leggere elemento (riga 1, colonna 2)
(aref *m* 1 2)           ; => 0

;; Scrivere
(setf (aref *m* 1 1) 5)
(aref *m* 1 1)           ; => 5

;; Dimensioni
(array-dimensions *m*)   ; => (3 3)
(array-dimension  *m* 0) ; => 3  (numero di righe)
(array-dimension  *m* 1) ; => 3  (numero di colonne)

;; Popolare con valori: matrice identita'
(defun identita (n)
  (let ((m (make-array (list n n) :initial-element 0)))
    (dotimes (i n m)
      (setf (aref m i i) 1))))

(identita 3)
; => #2A((1 0 0) (0 1 0) (0 0 1))
\end{lstlisting}
\end{esempio}

\begin{confronto}
\textbf{Array vs lista di liste.}
Una matrice si può rappresentare anche come lista di liste:
\texttt{((1 0 0) (0 1 0) (0 0 1))}.
Con la lista di liste, \fn{(nth j (nth i m))} ha costo $O(i+j)$.
Con l'array 2D, \fn{(aref m i j)} ha costo $O(1)$.
Per algoritmi su matrici grandi (trasposta, moltiplicazione,
algoritmi di grafo su matrice di adiacenza) usa sempre l'array.
\end{confronto}

\begin{workbook}
\textbf{Esercizi 101--104}
(Matrice, trasposta, inversione array, moltiplicazione matrici)\\[4pt]
L'esercizio~101 è l'essenziale: crea una matrice, leggi e scrivi celle.
L'esercizio~103 (inversione in-place) mostra la differenza tra
stile funzionale (crea un nuovo array) e stile imperativo
(modifica l'esistente): entrambi sono validi in Common Lisp.
\end{workbook}

% ================================================================
\section{\texttt{defstruct} --- strutture con campi nominati}
% ================================================================

\begin{concetto}{\texttt{defstruct}}
\keyline{}{Definisce un tipo con campi nominati, costruttore e accessor\\
generati automaticamente.}
\end{concetto}

\begin{tcolorbox}[
  enhanced,
  colback=black!2, colframe=black!20,
  arc=3pt, boxrule=0.4pt,
  left=12pt, right=12pt, top=9pt, bottom=9pt,
  before skip=8pt, after skip=10pt
]
\small\color{black!65}

\textbf{Cosa genera \fn{defstruct} automaticamente.}

\medskip
Quando scrivi \fn{(defstruct persona nome eta citta)}, Common Lisp
genera in automatico:
\begin{itemize}
  \item \fn{make-persona} --- costruttore con parametri keyword
  \item \fn{persona-nome}, \fn{persona-eta}, \fn{persona-citta} --- accessor (getter)
  \item \fn{persona-p} --- predicato di tipo
  \item \fn{copy-persona} --- copia superficiale della struttura
  \item Supporto per \fn{setf} su ogni accessor (setter)
\end{itemize}

\medskip
Il valore di ritorno è un oggetto \fn{\#S(PERSONA ...)},
che può essere stampato e riletto dal reader Lisp.
A differenza di una alist o hash table, il tipo è verificabile a runtime.
\end{tcolorbox}

\begin{esempio}[\fn{defstruct}: definizione e uso]
\begin{lstlisting}
;; Definizione: campi con valori di default
(defstruct persona
  nome
  (eta 18)         ; default: 18
  (citta "Roma"))  ; default: "Roma"

;; Costruzione con make-<nome>
(defparameter *p*
  (make-persona :nome "Mario" :eta 23 :citta "Milano"))

;; Accesso con accessor
(persona-nome *p*)    ; => "Mario"
(persona-eta  *p*)    ; => 23

;; Modifica con setf
(setf (persona-eta *p*) 24)
(persona-eta *p*)     ; => 24

;; Predicato di tipo
(persona-p *p*)       ; => T
(persona-p '(1 2))    ; => NIL

;; Valori di default
(make-persona :nome "Anna")
; => #S(PERSONA :NOME "Anna" :ETA 18 :CITTA "Roma")
\end{lstlisting}
\end{esempio}

\begin{esempio}[Funzioni su strutture]
\begin{lstlisting}
;; Validazione nel costruttore
(defun crea-persona (nome eta citta)
  (unless (> eta 0)
    (error "eta' deve essere positiva: ~a" eta))
  (make-persona :nome nome :eta eta :citta citta))

;; Stampa personalizzata
(defun stampa-persona (p)
  (format t "~a, ~a anni, ~a~%"
    (persona-nome p)
    (persona-eta  p)
    (persona-citta p)))

;; Sort su strutture
(sort (list (make-persona :nome "B" :eta 30)
            (make-persona :nome "A" :eta 25)
            (make-persona :nome "C" :eta 20))
      #'<
      :key #'persona-eta)
; ordinate per eta' crescente
\end{lstlisting}
\end{esempio}

\begin{workbook}
\textbf{Esercizio 105}
(Struttura \fn{persona} con validazione)\\[4pt]
Questo esercizio introduce \fn{defstruct}.
Dopo averlo fatto, nota come il codice diventa più leggibile:
\fn{(persona-nome p)} è più chiaro di \fn{(cdr (assoc 'nome p))}.
\end{workbook}

% ================================================================
\section{Alberi binari e BST}
% ================================================================

\begin{concetto}{Albero binario come lista}
\keyline{}{Convenzione: \fn{(valore sinistra destra)}.}
\keyline{\fn{nil}}{albero vuoto (foglia assente).}
\end{concetto}

\begin{esempio}[Albero binario: costruzione e accesso]
\begin{lstlisting}
;; Convenzione: (valore sottoalbero-sx sottoalbero-dx)
;; nil = albero vuoto

;; Nodi foglia
(defun foglia (v) (list v nil nil))

;; Accessor
(defun valore-nodo (n) (car   n))
(defun figlio-sx   (n) (cadr  n))
(defun figlio-dx   (n) (caddr n))

;; Costruire l'albero:
;;        5
;;       / \
;;      3   8
;;     / \
;;    1   4
(defparameter *albero*
  (list 5
    (list 3 (foglia 1) (foglia 4))
    (foglia 8)))

;; Accedere
(valore-nodo *albero*)            ; => 5
(valore-nodo (figlio-sx *albero*)) ; => 3
\end{lstlisting}
\end{esempio}

\argomento{Binary Search Tree (BST)}

\begin{concetto}{BST: la proprieta' fondamentale}
\keyline{}{Tutti i nodi nel sottoalbero sinistro sono \emph{minori}\\
del nodo corrente; tutti quelli a destra sono \emph{maggiori}.}

\boxsep
\detail{Questa proprietà garantisce ricerca, inserimento e
cancellazione in $O(\log n)$ su un albero bilanciato.
In Lisp si implementa in stile funzionale: ogni operazione restituisce
un nuovo albero, senza modificare quello originale.}
\end{concetto}

\begin{esempio}[BST: inserimento e ricerca funzionale]
\begin{lstlisting}
;; Inserimento: restituisce un NUOVO albero
(defun bst-insert (val tree)
  (cond ((null tree) (foglia val))         ; albero vuoto: crea foglia
        ((= val (valore-nodo tree)) tree)  ; gia' presente
        ((< val (valore-nodo tree))        ; va a sinistra
         (list (valore-nodo tree)
               (bst-insert val (figlio-sx tree))
               (figlio-dx tree)))
        (t                                 ; va a destra
         (list (valore-nodo tree)
               (figlio-sx tree)
               (bst-insert val (figlio-dx tree))))))

;; Ricerca
(defun bst-find (val tree)
  (cond ((null tree) nil)
        ((= val (valore-nodo tree)) t)
        ((< val (valore-nodo tree))
         (bst-find val (figlio-sx tree)))
        (t
         (bst-find val (figlio-dx tree)))))

;; Costruire un BST da una lista
(defun lista->bst (lst)
  (reduce #'(lambda (tree val) (bst-insert val tree))
          lst :initial-value nil))

(defparameter *bst* (lista->bst '(5 3 8 1 4)))
(bst-find 4 *bst*)   ; => T
(bst-find 7 *bst*)   ; => NIL
\end{lstlisting}
\end{esempio}

\begin{esempio}[Visita inorder: produce sequenza ordinata]
\begin{lstlisting}
;; Inorder (sinistra -> radice -> destra) su BST
;; produce sempre la sequenza ordinata crescente
(defun inorder (tree)
  (if (null tree)
      nil
      (append (inorder (figlio-sx tree))
              (list (valore-nodo tree))
              (inorder (figlio-dx tree)))))

(inorder *bst*)   ; => (1 3 4 5 8)  ordinata!
\end{lstlisting}
\end{esempio}

\begin{workbook}
\textbf{Esercizi 105--112}
(\fn{defstruct}, rappresentazione albero, BST, traversal,
altezza, specchia, bilanciato)\\[4pt]
Priorità: \textbf{106, 108, 109}.
L'esercizio~109 (traversal) mostra perché inorder su BST produce
una sequenza ordinata: è una proprietà non ovvia che vale la pena
verificare sperimentalmente.\\[4pt]
L'esercizio~112 (bilanciato) è il più difficile del gruppo:
leggi il suggerimento prima di tentare.
\end{workbook}

% ================================================================
\section{Pila e coda come ADT}
% ================================================================

\begin{concetto}{Abstract Data Type (ADT)}
\keyline{}{Un ADT espone solo le operazioni, non la rappresentazione interna.}

\boxsep
\detail{La pila e la coda si implementano su liste Lisp,
ma il chiamante non sa (e non deve sapere) come sono fatte dentro.
È incapsulamento: se si cambia la rappresentazione interna,
l'interfaccia non cambia.}
\end{concetto}

\begin{esempio}[Pila (stack) --- LIFO]
\begin{lstlisting}
;; Rappresentazione: una cons-cell (lista . puntatore-fine)
;; ma per la pila basta una lista: la testa e' il top

(defun make-stack () (list nil))        ; wrapper mutabile
(defun stack-empty-p (s) (null (car s)))
(defun stack-peek (s)    (caar s))

(defun stack-push (item s)
  (setf (car s) (cons item (car s))))

(defun stack-pop (s)
  (if (stack-empty-p s)
      (error "stack vuoto")
      (let ((top (caar s)))
        (setf (car s) (cdar s))
        top)))

;; Uso
(defparameter *s* (make-stack))
(stack-push 1 *s*)
(stack-push 2 *s*)
(stack-push 3 *s*)
(stack-pop *s*)    ; => 3
(stack-peek *s*)   ; => 2
\end{lstlisting}
\end{esempio}

\begin{esempio}[Applicazione: validazione parentesi]
\begin{lstlisting}
(defun valida-parentesi (stringa)
  (let ((s (make-stack))
        (aperte '(#\( #\[ #\{))
        (chiuse '(#\) #\] #\})))
    (every #'identity
      (loop for c across stringa
            collect
            (cond
              ((member c aperte) (stack-push c s) t)
              ((member c chiuse)
               (and (not (stack-empty-p s))
                    (char= (stack-pop s)
                           (nth (position c chiuse) aperte))))
              (t t)))
      (stack-empty-p s))))

(valida-parentesi "([]{})")   ; => T
(valida-parentesi "([)]")     ; => NIL
\end{lstlisting}
\end{esempio}

\begin{workbook}
\textbf{Esercizi 113--115}
(Pila, coda, validazione parentesi)\\[4pt]
L'esercizio~115 (validazione parentesi) è l'applicazione classica
della pila: usata in tutti i parser, compilatori, editor.
Dopo averlo risolto avrai un'intuizione concreta di perché
la pila è la struttura giusta per problemi di matching.
\end{workbook}

% ================================================================
\section*{Riepilogo del capitolo}
% ================================================================

\begin{tcolorbox}[
  enhanced,
  colback=csBlue!5, colframe=csBlue!40,
  arc=4pt, boxrule=0.6pt,
  left=14pt, right=14pt, top=12pt, bottom=12pt,
  before skip=16pt, after skip=16pt
]
\textbf{Quello che hai imparato:}

\medskip
\begin{itemize}
  \item \textbf{Alist}: lista di coppie \fn{(chiave . valore)}.
        \fn{assoc} per cercare, \fn{acons} per aggiungere.
        Funzionale, $O(n)$. Buona per dizionari piccoli.
  \item \textbf{Hash table}: \fn{make-hash-table}, \fn{gethash},
        \fn{setf gethash}, \fn{remhash}, \fn{maphash}.
        $O(1)$ per accesso. Per stringhe usa \fn{:test \#'equal}.
  \item \textbf{Array 2D}: \fn{make-array '(n m)}, \fn{aref} con
        due indici. Per matrici e algoritmi con accesso casuale.
  \item \textbf{\fn{defstruct}}: genera costruttore, accessor,
        predicato. Campi nominati, type-checking, \fn{setf} su accessor.
  \item \textbf{BST}: proprietà sinistra-minore/destra-maggiore.
        Inserimento e ricerca funzionali ($O(\log n)$).
        Inorder produce sequenza ordinata.
  \item \textbf{Pila/coda}: ADT con interfaccia esplicita.
        La rappresentazione è privata; l'interfaccia è stabile.
\end{itemize}

\medskip
\textbf{Prossimo:} Capitolo~6 --- \emph{Macro}.
Trasformare codice in codice prima della valutazione.
\end{tcolorbox}

\begin{workbook}
\textbf{Prima di andare al Capitolo~6:}\\[4pt]
Completa gli esercizi \textbf{93--120} (capitolo
\emph{Strutture dati}).\\[4pt]
Priorità: \textbf{93, 95, 101, 105, 106, 108, 113, 115}.
Questi otto coprono una struttura per tipo.\\[4pt]
Gli esercizi 117--120 (grafo, DFS, BFS, conta parole) sono
collaterali ma molto utili: applicano le strutture dati a
problemi reali.
\end{workbook}
